\section{Matrices}

A matrix is one of the first mathematical objects that looks very simple and very mechanical. It is a rectangular table of numbers. At first sight, this may suggest that matrix theory is only a matter of arranging numbers and following rules. This would be a serious misunderstanding. A matrix is useful because it allows us to store information, describe relations, encode systems of equations, and represent transformations. The same object can be read as a table, as a rule of calculation, as a linear transformation, or as a compact description of a structured problem.

In this section we will learn the basic operations on matrices: addition, subtraction, multiplication by a number, multiplication of two matrices, powers of a matrix, and simple operations on rows. The goal is not only to perform these operations correctly. The goal is to understand what kind of object is being manipulated and why the operation is defined in that particular way.

\subsection{What is a matrix?}

Before we calculate with matrices, we should learn to read them. A matrix has rows and columns. The first question is therefore not ``What formula should I use?'', but ``What is the size of the object in front of me?''

\begin{definitionbox}
A matrix with \(m\) rows and \(n\) columns is called an \(m \times n\) matrix. It has the form
\[
A=
\begin{pmatrix}
a_{11} & a_{12} & \cdots & a_{1n}\\
a_{21} & a_{22} & \cdots & a_{2n}\\
\vdots & \vdots & \ddots & \vdots\\
a_{m1} & a_{m2} & \cdots & a_{mn}
\end{pmatrix}.
\]
The number \(a_{ij}\) is the entry in row \(i\) and column \(j\).
\end{definitionbox}

The notation \(a_{ij}\) is not decoration. It tells us how to locate information. The first index tells us the row, and the second index tells us the column. Thus \(a_{23}\) means: second row, third column. This is a small convention, but conventions of this kind are what allow mathematics to be precise.

\begin{examplebox}
Consider the matrix
\[
A=
\begin{pmatrix}
1 & 2 & 5\\
0 & -1 & 4
\end{pmatrix}.
\]
This is a \(2\times 3\) matrix. It has two rows and three columns. The entry \(a_{12}\) is equal to \(2\), because it is in the first row and second column. The entry \(a_{23}\) is equal to \(4\), because it is in the second row and third column.
\end{examplebox}

A row can be read as a horizontal list of numbers. A column can be read as a vertical list of numbers. Later, when matrices are used to describe systems of equations or transformations, this distinction will become important. For now it is enough to remember that a matrix is not only a pile of numbers: it is a pile of numbers with a shape.

\subsection{Equality of matrices}

Two matrices can be equal only if they have the same size. After that, equality is checked entry by entry.

\begin{definitionbox}
Let \(A=(a_{ij})\) and \(B=(b_{ij})\) be matrices of the same size \(m\times n\). We say that \(A=B\) if
\[
a_{ij}=b_{ij}
\]
for every row index \(i\) and every column index \(j\).
\end{definitionbox}

This definition is simple, but it teaches an important habit. When comparing mathematical objects, we first ask whether they are objects of the same type. A number can be equal to a number, a vector to a vector of the same dimension, and a matrix to a matrix of the same size. If the shapes do not match, the comparison is already suspicious.

\subsection{Addition, subtraction, and multiplication by a scalar}

Addition of matrices is defined entry by entry. This is natural when two matrices represent the same kind of table. If two tables have the same shape, then the entry in position \((i,j)\) in the sum should come from adding the entries in position \((i,j)\) in the two original tables.

\begin{definitionbox}
Let \(A=(a_{ij})\) and \(B=(b_{ij})\) be matrices of the same size. Their sum is the matrix
\[
A+B=(a_{ij}+b_{ij}).
\]
Their difference is the matrix
\[
A-B=(a_{ij}-b_{ij}).
\]
If \(\lambda\) is a number, then
\[
\lambda A=(\lambda a_{ij}).
\]
\end{definitionbox}

The condition ``of the same size'' is essential. We cannot add a \(2\times 3\) matrix to a \(3\times 2\) matrix, because there is no consistent way to pair the entries position by position.

\begin{examplebox}
Let
\[
A=\begin{pmatrix}1 & 2\\ 3 & 4\end{pmatrix},
\qquad
B=\begin{pmatrix}0 & -1\\ 2 & 1\end{pmatrix}.
\]
Then
\[
A+B=
\begin{pmatrix}
1+0 & 2+(-1)\\
3+2 & 4+1
\end{pmatrix}
=
\begin{pmatrix}
1 & 1\\
5 & 5
\end{pmatrix}.
\]
Similarly,
\[
A-B=
\begin{pmatrix}
1-0 & 2-(-1)\\
3-2 & 4-1
\end{pmatrix}
=
\begin{pmatrix}
1 & 3\\
1 & 3
\end{pmatrix},
\]
and
\[
2A=
\begin{pmatrix}
2 & 4\\
6 & 8
\end{pmatrix}.
\]
For a slightly longer expression, we calculate each part carefully:
\[
3B=\begin{pmatrix}0 & -3\\ 6 & 3\end{pmatrix},
\qquad
2A=\begin{pmatrix}2 & 4\\ 6 & 8\end{pmatrix},
\]
so
\[
3B-2A=
\begin{pmatrix}0 & -3\\ 6 & 3\end{pmatrix}
-
\begin{pmatrix}2 & 4\\ 6 & 8\end{pmatrix}
=
\begin{pmatrix}-2 & -7\\ 0 & -5\end{pmatrix}.
\]
\end{examplebox}

The important point is not only the final matrices. The important point is the method: we keep the shape, and we work position by position.

\subsection{Matrix multiplication: row meets column}

Matrix multiplication is different. It is not defined entry by entry. The reason is that multiplication is meant to combine rows of the first matrix with columns of the second matrix. This operation is less obvious than addition, but it is one of the main reasons matrices are powerful.

\begin{definitionbox}
Let \(A=(a_{ij})\) be an \(m\times n\) matrix and let \(B=(b_{ij})\) be an \(n\times p\) matrix. The product \(AB\) is the \(m\times p\) matrix \(C=(c_{ij})\), where
\[
c_{ij}=a_{i1}b_{1j}+a_{i2}b_{2j}+\cdots+a_{in}b_{nj}.
\]
In words: the entry \(c_{ij}\) is obtained by multiplying row \(i\) of \(A\) by column \(j\) of \(B\).
\end{definitionbox}

The sizes tell us whether multiplication is possible:
\[
(m\times n)(n\times p)=m\times p.
\]
The inner dimensions must match. The outer dimensions give the size of the answer.

\begin{remarkbox}
A useful way to read matrix multiplication is this: each entry of the product is a dot product of one row and one column. Therefore matrix multiplication is not just a bigger version of ordinary multiplication. It is a rule for combining structured information.
\end{remarkbox}

\begin{examplebox}
Using the same matrices
\[
A=\begin{pmatrix}1 & 2\\ 3 & 4\end{pmatrix},
\qquad
B=\begin{pmatrix}0 & -1\\ 2 & 1\end{pmatrix},
\]
we calculate
\[
AB=
\begin{pmatrix}
1\cdot 0+2\cdot 2 & 1\cdot (-1)+2\cdot 1\\
3\cdot 0+4\cdot 2 & 3\cdot (-1)+4\cdot 1
\end{pmatrix}
=
\begin{pmatrix}
4 & 1\\
8 & 1
\end{pmatrix}.
\]
Every entry has a small story. The entry in the first row and first column comes from the first row of \(A\) and the first column of \(B\):
\[
(1,2)\cdot (0,2)=1\cdot 0+2\cdot 2=4.
\]
The entry in the second row and first column comes from
\[
(3,4)\cdot (0,2)=3\cdot 0+4\cdot 2=8.
\]
\end{examplebox}

This is the kind of calculation in which one should not rush. The entries of the answer are not found by matching positions as in addition. They are found by moving through rows and columns.

\subsection{The order of multiplication matters}

For numbers, multiplication is commutative: \(ab=ba\). For matrices, this is usually false. Sometimes both products \(AB\) and \(BA\) are defined, but they may give different answers.

\begin{examplebox}
For the matrices from the previous example, we already found
\[
AB=
\begin{pmatrix}
4 & 1\\
8 & 1
\end{pmatrix}.
\]
Now calculate the opposite product:
\[
BA=
\begin{pmatrix}0 & -1\\ 2 & 1\end{pmatrix}
\begin{pmatrix}1 & 2\\ 3 & 4\end{pmatrix}
=
\begin{pmatrix}
-3 & -4\\
5 & 8
\end{pmatrix}.
\]
Thus
\[
AB\ne BA.
\]
The two products use the same two matrices, but they ask different row-column questions.
\end{examplebox}

This example is not a small technical exception. It reveals a central feature of matrix algebra: order carries meaning. If matrices represent operations, then \(AB\) means that one operation is performed after another. Changing the order may change the result.

\begin{theorembox}
For matrices of compatible sizes, the following laws hold:
\[
A+B=B+A,
\]
\[
(A+B)+C=A+(B+C),
\]
\[
(AB)C=A(BC),
\]
\[
A(B+C)=AB+AC,
\qquad
(A+B)C=AC+BC.
\]
However, in general,
\[
AB\ne BA.
\]
\end{theorembox}

The last statement should be remembered carefully. It does not say that \(AB\) is never equal to \(BA\). It says that equality cannot be assumed. If we want to claim that two matrices commute, we must check it or prove it from special structure.

\subsection{Zero, identity, and diagonal matrices}

Some matrices play special roles. The zero matrix has all entries equal to zero. The identity matrix is the matrix that does nothing under multiplication, in the same way that the number \(1\) does nothing under ordinary multiplication.

\begin{definitionbox}
The \(n\times n\) identity matrix is
\[
I_n=
\begin{pmatrix}
1 & 0 & \cdots & 0\\
0 & 1 & \cdots & 0\\
\vdots & \vdots & \ddots & \vdots\\
0 & 0 & \cdots & 1
\end{pmatrix}.
\]
If \(A\) has a compatible size, then
\[
I A=A,
\qquad
A I=A.
\]
A diagonal matrix is a square matrix whose entries outside the main diagonal are zero.
\end{definitionbox}

Diagonal matrices are especially easy to understand. They act independently on each coordinate. This is why they are often the first matrices in which one can see structure instead of only arithmetic.

\begin{examplebox}
Let
\[
D=\operatorname{diag}(2,-3,5)
=
\begin{pmatrix}
2 & 0 & 0\\
0 & -3 & 0\\
0 & 0 & 5
\end{pmatrix}
\]
and let
\[
E=\operatorname{diag}(a,b,c)
=
\begin{pmatrix}
a & 0 & 0\\
0 & b & 0\\
0 & 0 & c
\end{pmatrix}.
\]
Then
\[
DE=\operatorname{diag}(2a,-3b,5c)
\]
and
\[
ED=\operatorname{diag}(2a,-3b,5c).
\]
Therefore \(DE=ED\). In this special case the matrices commute because every coordinate is scaled independently, and no coordinate is mixed with another one.
\end{examplebox}

Powers of a diagonal matrix are also easy:
\[
D^3=\operatorname{diag}(2^3,(-3)^3,5^3)
=\operatorname{diag}(8,-27,125).
\]
If no diagonal entry is zero, then the inverse diagonal matrix is obtained by taking reciprocals:
\[
D^{-1}=\operatorname{diag}\left(\frac12,-\frac13,\frac15\right).
\]
A complete study of inverse matrices will come later. Here the diagonal case is useful because the idea is visible without complicated computation.

\subsection{Rows, columns, and elementary row operations}

Matrices are often manipulated row by row. A row operation changes the matrix, but it does so in a controlled and visible way. Such operations will later be used to solve systems of linear equations.

\begin{examplebox}
Consider
\[
C=
\begin{pmatrix}
1 & 0 & 2\\
-1 & 3 & 1\\
0 & 2 & -1
\end{pmatrix}.
\]
First swap the first and third rows. We obtain
\[
C_1=
\begin{pmatrix}
0 & 2 & -1\\
-1 & 3 & 1\\
1 & 0 & 2
\end{pmatrix}.
\]
Now add twice the new first row to the second row. The second row becomes
\[
(-1,3,1)+2(0,2,-1)=(-1,7,-1).
\]
Therefore the final matrix is
\[
C_2=
\begin{pmatrix}
0 & 2 & -1\\
-1 & 7 & -1\\
1 & 0 & 2
\end{pmatrix}.
\]
\end{examplebox}

The point of writing the intermediate step is not to make the solution longer. It makes the reasoning inspectable. We can see exactly which row changed and why.

\subsection{Column vectors as matrices}

A vector can be written as a matrix with one column. This is not just a change of notation. It allows us to use matrix operations to organize vector calculations.

Let
\[
u=\begin{pmatrix}1\\ -2\\ 3\end{pmatrix},
\qquad
v=\begin{pmatrix}2\\ 0\\ -1\end{pmatrix}.
\]
Then
\[
u+v=
\begin{pmatrix}3\\ -2\\ 2\end{pmatrix},
\qquad
u-v=
\begin{pmatrix}-1\\ -2\\ 4\end{pmatrix}.
\]
The transpose of a column vector is a row vector:
\[
v^{\top}=\begin{pmatrix}2 & 0 & -1\end{pmatrix}.
\]
Thus the product \(uv^{\top}\) is a \(3\times 3\) matrix:
\[
uv^{\top}
=
\begin{pmatrix}1\\ -2\\ 3\end{pmatrix}
\begin{pmatrix}2 & 0 & -1\end{pmatrix}
=
\begin{pmatrix}
2 & 0 & -1\\
-4 & 0 & 2\\
6 & 0 & -3
\end{pmatrix}.
\]
Similarly,
\[
vu^{\top}
=
\begin{pmatrix}2\\ 0\\ -1\end{pmatrix}
\begin{pmatrix}1 & -2 & 3\end{pmatrix}
=
\begin{pmatrix}
2 & -4 & 6\\
0 & 0 & 0\\
-1 & 2 & -3
\end{pmatrix}.
\]
These two matrices are not the same. Again, the order of multiplication matters.

There is also a structural observation hidden in this example. In the matrix \(uv^{\top}\), every column is a multiple of \(u\). Therefore all columns point in the same direction in the sense of linear dependence. This is why, if \(u\) and \(v\) are non-zero, the matrix \(uv^{\top}\) has rank \(1\). We do not need the full theory of rank yet, but the example shows how a calculation can reveal structure.

\subsection{Powers of a matrix}

Powers of a matrix are defined only for square matrices. If \(P\) is square, then
\[
P^2=PP,
\qquad
P^3=PPP,
\]
and so on. Matrix powers are useful when the same operation is repeated.

\begin{examplebox}
Let
\[
P=\begin{pmatrix}
1 & 1 & 0\\
0 & 1 & 1\\
1 & 0 & 1
\end{pmatrix}.
\]
Then
\[
P^2=
\begin{pmatrix}
1 & 2 & 1\\
1 & 1 & 2\\
2 & 1 & 1
\end{pmatrix}
\]
and
\[
P^3=
\begin{pmatrix}
2 & 3 & 3\\
3 & 2 & 3\\
3 & 3 & 2
\end{pmatrix}.
\]
The computation is ordinary matrix multiplication, but after calculating we should pause. The matrices show a cyclic pattern: the rows are related by rotation of positions. This is not an accident; it reflects the structure already present in \(P\).
\end{examplebox}

A good exercise is not finished when the arithmetic ends. It is finished when we have also asked what the arithmetic reveals.

\subsection{A first glimpse: matrices as transformations}

One reason matrices become powerful is that they can describe transformations. For example, the matrix
\[
R(\theta)=
\begin{pmatrix}
\cos\theta & -\sin\theta\\
\sin\theta & \cos\theta
\end{pmatrix}
\]
represents rotation of the plane by angle \(\theta\) around the origin. This interpretation is not needed to multiply the matrices, but it helps explain why the product should have a meaningful form.

If we rotate by \(\theta_2\) and then by \(\theta_1\), the total effect should be a rotation by \(\theta_1+\theta_2\). Algebra confirms this:
\[
R(\theta_1)R(\theta_2)=
\begin{pmatrix}
\cos\theta_1 & -\sin\theta_1\\
\sin\theta_1 & \cos\theta_1
\end{pmatrix}
\begin{pmatrix}
\cos\theta_2 & -\sin\theta_2\\
\sin\theta_2 & \cos\theta_2
\end{pmatrix}.
\]
The first entry of the product is
\[
\cos\theta_1\cos\theta_2-\sin\theta_1\sin\theta_2
=\cos(\theta_1+\theta_2).
\]
The other entries are obtained in the same way from trigonometric addition formulas, and the result is
\[
R(\theta_1)R(\theta_2)=R(\theta_1+\theta_2).
\]
This is a beautiful moment: a matrix calculation, a geometric transformation, and a trigonometric identity all describe the same structure.

\subsection{How to work with matrix problems}

When working with matrices, do not begin by searching for a memorized rule. Begin by reading the object.

\begin{summarybox}
For matrix problems, ask the following questions.
\begin{itemize}
\item What is the size of each matrix?
\item Are the operations even defined?
\item If I add or subtract matrices, am I working entry by entry?
\item If I multiply matrices, which row meets which column?
\item Does the order of multiplication matter in this problem?
\item Is there special structure: zero matrix, identity matrix, diagonal matrix, symmetry, repeated rows, or a visible pattern?
\item After the computation, what did the result tell me?
\end{itemize}
\end{summarybox}

The last question is essential. Matrix algebra is not only a technique for producing new arrays of numbers. It is a language for seeing structure. A student who learns to pause, read the objects, and interpret the result is already doing more than mechanical calculation.
